\section{Introduction}
Let $H$ be a complex Hilbert space and let $\B(H)$ denote its bounded
linear operators. The \emph{commutant} of $S\in\B(H)$ is
\[
 \{S\}'=\{A\in\B(H):AS=SA\}.
\]
A closed subspace is \emph{hyperinvariant} for $S$ if every operator in
$\{S\}'$ leaves it invariant. An operator algebra is \emph{transitive}
if its only closed common invariant subspaces are $\{0\}$ and $H$.
Thus $S$ has no nonzero proper closed hyperinvariant subspace precisely
when its commutant is transitive.

The hyperinvariant-subspace problem asks whether every bounded nonscalar
operator on an infinite-dimensional separable complex Hilbert space has
a nonzero proper closed hyperinvariant subspace. The construction below
gives a negative answer. Its purpose is to exhibit the mechanism
explicitly: the ordinary invariant subspaces of each weighted-shift
fiber are coordinate tails, and a second, commuting operator moves those
tails backward. Invariance under the commutant therefore imposes a
strict descent that cannot persist on a probability space.

\subsection{The construction theorem}
Put $K=\ell^2(\Z)$, with orthonormal basis $(e_j)_{j\in\Z}$, and define
\[
 \Ktail{h}=\overline{\operatorname{span}}\{e_j:j\ge h\}
 \quad(h\in\Z),\qquad
 \Ktail{-\infty}=K,\qquad \Ktail{+\infty}=\{0\}.
\]
Let $X=\Z_2$ be the compact additive group of $2$-adic integers, the
inverse limit of the groups $\Z/2^m\Z$. Equip it with completed Haar
probability $\mu$, put $\theta x=x+1$, and set $H_0=L^2(X,\mu;K)$.
For $g\in L^\infty(X)$, the scalar multiplier is
$(M_gf)(x)=g(x)f(x)$.

\begin{theorem}\label{thm:direct-model}
There are measurable positive weights $\beta_j:X\to(0,\infty)$,
$j\in\Z$, and a contraction $V\in\B(H_0)$ with the following properties.
The shifts
\[
 S_xe_j=\beta_j(x)e_{j+1},\qquad (Sf)(x)=S_xf(x),
\]
define a bounded operator $S$ on $H_0$, and:
\begin{enumerate}
\item For every $x\in X$, the closed invariant subspaces of $S_x$ are
exactly $\Ktail{h}$, $h\in\Z\cup\{-\infty,+\infty\}$.
\item The operator $S$ is nonzero, $\|S\|=e^{-10}$, and for every
integer $n\ge1$,
\begin{equation}\label{eq:main-power-bound}
 \|S^n\|\le
 \exp\left(-10\left\lfloor\frac{(n+1)^2}{4}\right\rfloor\right).
\end{equation}
\item The operator $V$ commutes with $S$ and has the form
\[
 (Vf)(\theta x)=V_xf(x),\qquad V_xe_j=b_j(x)e_{j-1},
\]
where $0\le b_j(x)\le1$ are measurable and are positive on a common
conull set. On that set,
\[
 \overline{V_x\Ktail{h}}=\Ktail{h-1}\quad(h\in\Z),
 \qquad \overline{V_xK}=K.
\]
\item The unital complex algebra generated by $S$, $V$, and all
scalar multipliers $M_g$ is transitive.
\end{enumerate}
\end{theorem}

All generators in the last assertion commute with $S$. The explicit
power bound implies norm-quasinilpotence,
$\lim_{n\to\infty}\|S^n\|^{1/n}=0$, while the first assertion will give
many ordinary invariant subspaces of the resulting operator.

\begin{corollary}\label{cor:direct-algebra}
Every infinite-dimensional separable complex Hilbert space admits a
nonzero bounded norm-quasinilpotent operator with no nonzero proper
closed hyperinvariant subspace. Its commutant is a proper strongly
closed unital transitive complex operator algebra.
\end{corollary}

Here the strong operator topology is pointwise norm convergence on
$H$. Consequently the commutant in the corollary also gives a negative
answer to the question whether every unital transitive operator algebra
is strongly dense in $\B(H)$.

\subsection{Context and relation to the companion}
Douglas and Pearcy studied the connection between hyperinvariant
subspaces and transitive algebras~\cite{DouglasPearcy1972}. A major
positive result is Arveson's density theorem: a transitive algebra on
a separable Hilbert space which contains a maximal abelian self-adjoint
subalgebra of $\B(H)$ is strongly dense in $\B(H)$
\cite[Theorem~3.3]{Arveson1967}. The scalar multiplier algebra used here
has infinite multiplicity. Section~\ref{sec:structure} gives an explicit
operator commuting with every such multiplier that is not itself a
scalar multiplier, so the maximality hypothesis is absent.

Domar's work on invariant subspaces of weighted translation spaces
\cite{Domar1981} and Grabiner's work on unicellular shifts
\cite{Grabiner1982} provide the classical lattice ideas behind the
fibers. Section~\ref{sec:criterion} proves the precise quadratic
weighted-translation criterion we need. Its convex-envelope and
support-separation argument follows Domar's method. The proof includes
both normalized convolution-tail estimates; no general criterion for
arbitrary weights is assumed. Our additional task is to link the fibers
by a bounded commuting operator that moves their invariant tails
backward. Opposite logarithmic changes at the two ends of the bilateral
sequence make this possible, and one measurable normalization controls
the operator over the whole base.

Atzmon and Sodin constructed completely indecomposable bilateral
weighted shifts with spectrum the unit circle~\cite{AtzmonSodin1999};
their 1999 paper left open whether those shifts have proper hyperinvariant
subspaces. Here the extra commutant action is supplied by transport
between different fibers. For a study of operators with shift-type invariant subspaces,
see K\'erchy~\cite{Kerchy2026}.
The invariant- and hyperinvariant-subspace problems must remain
distinct: the K\'erchy--Pearcy preprint claiming an equivalence was
withdrawn after a counterexample to its Theorem~3
\cite{KerchyPearcy2025Withdrawn}. The operator constructed here preserves
every fixed tail $L^2(X;\Ktail{h})$; its hyperinvariant conclusion is
proved directly.

The companion paper \cite{CircleCompanion} constructs, for every
prescribed irrational rotation, a continuous circle weight whose
operator in the rotation factor has no nontrivial invariant projection
in that factor. It derives the existence conclusion of
Corollary~\ref{cor:direct-algebra} by the usual projection argument;
see also Dykema~\cite[Proposition~2.1]{Dykema2005}.
The present construction has a different purpose: it exhibits the
ordinary invariant tails, the bounded operator that moves them backward,
and the explicit estimate~\eqref{eq:main-power-bound}. The two proofs
use different constructions. Our only proof imports from the companion
are the two elementary facts stated below.

\subsection{Proof overview and elementary inputs}
The proof first makes the possible fiber subspaces completely explicit.
Section~\ref{sec:criterion} classifies invariant subspaces of a weighted
translation whose logarithmic weights have opposite quadratic growth
at the two ends. Section~\ref{sec:fibers} constructs weights over
$\Z_2$ satisfying that criterion in every fiber. Their decay also gives
the uniform power estimate, which proves norm-quasinilpotence of the
operator on $H_0$ rather than merely of its individual fibers.

To couple the fibers, we prescribe a map from $e_j$ at $x$ to a positive
multiple of $e_{j-1}$ at $\theta x$. Commutation with the forward shifts
determines a recurrence for those multiples. Section~\ref{sec:intertwiner}
solves the recurrence in both directions. The logarithmic changes of the
weights under $\theta$ have opposite signs at the two ends; the
coefficients therefore decay in both directions. Dividing all coefficients in each fiber by their common
supremum yields the contraction $V$ without changing commutation.

Finally, a subspace invariant under every scalar multiplier has an
orthogonal projection represented by a measurable projection field.
Invariance under $S$ makes each fiber a tail $\Ktail{s(x)}$.
Section~\ref{sec:subspaces} uses the dense images under $V_x$ to obtain
\[
 s(\theta x)\le s(x)-1\qquad\text{whenever }s(x)\in\Z.
\]
Probability preservation excludes finite labels, and ergodicity excludes
a nonconstant mixture of full and zero fibers. Section~\ref{sec:structure}
then proves the corollary and records the structure of the example.

We use the following two elementary inputs, with the stated scope.
Their complete proofs are in the companion \cite{CircleCompanion}.
\begin{enumerate}
\item \emph{Fields from scalar multipliers}
\cite[\FieldImportLocator]{CircleCompanion}.
For a probability space $(Y,\nu)$, $K=\ell^2(\Z)$, and
$A\in\B(L^2(Y;K))$ commuting with every scalar multiplier,
there is an essentially bounded field $A_y\in\B(K)$, unique almost
everywhere, such that
\[
 (Af)(y)=A_yf(y),\qquad \|A_y\|\le\|A\|\quad\text{a.e.}
\]
For every $v\in K$, the map $y\mapsto A_yv$ is measurable.
Products and adjoints act fiberwise; an orthogonal projection has
orthogonal projection fibers outside one null set.
\item \emph{The commutant algebra}
\cite[\CommutantImportLocator]{CircleCompanion}.
For any bounded nonscalar operator on a complex Hilbert space, its
commutant is a proper strongly closed unital complex algebra. If its
only closed common invariant subspaces are zero and the whole space,
that algebra is transitive.
\end{enumerate}
The first input requires no base transformation, covariance, trace or
factor assumption. We apply it only after proving reduction by all
scalar multipliers. The second requires neither separability nor
quasinilpotence; the constructed operator will satisfy its nonscalarity
hypothesis. The field input is used in Section~\ref{sec:subspaces}, and the
commutant input in Section~\ref{sec:structure}.
